\section{Результати}
\label{sec:results}

У цьому розділі наведено результати чотирьох серій, описаних у Розділі~\ref{sec:design}.
Для ланцюгів використано нові запуски TwoScale після виправлення залишкового локального годинника та роботи з генератором випадкових чисел. Перераховано всі 48 конфігурацій, по 400 пар у кожній; еталонні результати Full залишилися незмінними. Рисунки й числові оцінки для ланцюгів побудовано за цими 19\,200 парами.
Для тестів з одиничним експортом наведено результати попередніх серій: вони не входили до цього перерахунку. Тому ці два тести характеризують попередню реалізацію TwoScale і не є перевіркою її виправленої версії.

Для двох запусків Full використано збережені продовження до 40\,000, щоб отримати спостережені значення $T_9$, як зазначено в Розділі~\ref{sec:design}. На спільному горизонті 30\,000 останньої спільноти досягли 19\,198 із 19\,200 запусків Full і всі 19\,200 запусків TwoScale. Подальші оцінки розподілів $T_9$ включають ці два продовження; частки завершення при різних горизонтах не порівнюємо.

\subsection{Одиничний експорт із кліки}
\label{sec:res_clique}

\emph{Часи експорту з джерела, яке швидко насичується, добре узгоджуються, але агрегований інтерфейс допускає неможливі порядки передачі.} У тесті «кліка--листок» (рисунок~\ref{fig:clique_leaf}) середні часи експорту становили 110.18 (\Full{}) і 112.62 (\TS{}); парна різниця $+2.43$ мала бутстреп-інтервал $[-0.35,5.27]$. Емпіричні функції розподілу майже збігаються ($W_1=3.25$, тобто $\eW=2.95\%$; відстань Колмогорова--Смирнова 0.0142, асимптотичне $p\approx0.27$). Із середнього часу експорту близько 110 лише $\mathbb{E}[T_B]=8.36$ у середньому потрібно, щоб заразити перший вузол-мост (аналітична оцінка для кліки, Додаток~\ref*{app:clique}); решта припадає на очікування передачі через інтерфейс з інтенсивністю не більшою за $k\beta=0.01$. Отже, період, протягом якого стани межі й усієї спільноти відрізняються, короткий, що пояснює малу спостережену різницю часів експорту в цьому тесті. Твердження~\ref{prop:ordering} описує напрям зміщення для ідеального процесу, але не встановлює величини похибки сіткової реалізації.

Водночас діагностика порядку розрізняє моделі: у 98 із 10\,000 запусків \TS{} (0.98\%) листок було заражено раніше за обидва вузли-мости, чого ніколи не трапляється у \Full{} (рисунок~\ref{fig:clique_leaf}b,c). На середній панелі показано $T_{\mathrm{export}}-T_B$: фізичний час очікування у \Full{}, але знаковий діагностичний показник порядку у \TS{}.

\begin{figure}[tbp]
\centering
\includegraphics[width=\linewidth]{Images/clique_leaf/clique_leaf_endpoint_and_mechanism.png}
\caption{Тест «кліка--листок» для попередньої реалізації TwoScale ($n=100$, $k=2$, $\beta=0.005$, 10\,000 запусків для кожної моделі). (a) Емпіричні функції розподілу часу експорту. (b) Емпіричні функції розподілу $T_{\mathrm{export}}-T_B$; від'ємні значення (заштрихована область) відповідають експорту до зараження будь-якого вузла-моста. (c) Частка таких запусків з обходом мостів із 95\% інтервалами. ``Full Micro'' позначає \Full{}, а ``MicroMacro'' --- \TS{}.}
\label{fig:clique_leaf}
\end{figure}

\subsection{Одиничний експорт зі зв'язних випадкових джерел}
\label{sec:res_source}

\emph{Зі збільшенням довжини внутрішніх шляхів обхід мостів стає поширеним, а агрегований експорт відбувається раніше, що за напрямом узгоджується з твердженням~\ref{prop:ordering} для ідеального процесу.} На рисунку~\ref{fig:source_sweep} показано серію експериментів «джерело--листок». Середні часи активації мостів у двох моделях були подібними: різниця не перевищувала 7\% для кожного ASPL (наприклад, 1165 і 1136 за ASPL 6, $k=1$), як і очікувалося через однакову локальну динаміку до експорту. У \Full{} за $k=1$ очікування після активації моста залишалося в межах 322--359 за кожного ASPL, що узгоджується з $\mathrm{Exp}(\beta)$ із середнім 333. У \TS{} знаковий інтервал $T_{\mathrm{export}}-T_B$ ($k=1$) зменшився приблизно з 340 у щільних джерелах до 207--295 за ASPL $\ge5$, оскільки агрегована інтенсивність накопичувалася ще до того, як справжній міст ставав інфекційним.

\begin{figure}[tbp]
\centering
\includegraphics[width=\linewidth]{Images/fig_aspl_export_fidelity.png}
\caption{Серія експериментів «джерело--листок» для попередньої реалізації TwoScale: $n=1000$, $\beta=0.003$, $k=1$ (угорі) та $k=5$ (унизу), 400 парних запусків для кожної точки. Ліворуч: середній час експорту. У центрі: відносна відстань Вассерштайна $\eW$ між розподілами часу експорту. Праворуч: частка експортів \TS{}, що передували зараженню першого моста; у \Full{} таких випадків не було. Відрізки похибок --- 95\% інтервали блокового бутстрепу. ``Full Micro'' позначає \Full{}, а ``MicroMacro'' --- \TS{}.}
\label{fig:source_sweep}
\end{figure}

Частка обходів зросла від не більш як 0.5\% за ASPL $\le1.75$ до 19.75\% (95\% інтервал 15.8--24.0\%) за ASPL 4 і $k=1$, а за ASPL 6 --- до 32.5\% для обох ширин інтерфейсу. Зсув середнього часу експорту був статистично відмінним від нуля, з інтервалами, що не містили нуль, для обох ширин за кожного ASPL $\ge5$. Його значення були в межах від $-6.2\%$ до $-12.4\%$; за ASPL 6 вони становили $-8.3\%$ ($[-13.1,-3.5]$) для $k=1$ і $-6.2\%$ ($[-10.3,-2.3]$) для $k=5$. Нижче ASPL 5 статистично відмінний від нуля зсув виявлено лише для $k=5$ за ASPL 3.75, 4.25 і 4.5 (від $-3.7\%$ до $-5.8\%$). У жодній точці серії не виявлено додатного зсуву, що узгоджується з напрямом, передбаченим твердженням~\ref{prop:ordering}.

Нерегулярну центральну панель слід тлумачити обережно. Для майже експоненційних часів експорту та $R=400$ запусків для кожної моделі очікуване $\eW$ двох вибірок з одного розподілу становить приблизно $\sqrt{\pi/R}\approx9\%$ (Додаток~\ref*{app:floor}). Отже, значення 3.7--8.6\%, спостережені за ASPL $\le1.75$, узгоджуються із суто вибірковою мінливістю. Тому відстані для одиничного експорту слід розглядати разом зі зсувом середнього та його інтервалом. Натомість обхід у \Full{} за побудовою дорівнює нулю, тому кожен спостережений обхід є справжньою розбіжністю. За ASPL 4 він стосується однієї п'ятої експортів, хоча розподіли часу експорту залишаються близькими.

\subsection{Повторні переходи через ланцюг клік}
\label{sec:res_cliques}

\emph{Коли внутрішнє поширення швидке, близьке узгодження зберігається після дев'яти послідовних агрегованих переходів.} У ланцюзі з десяти клік по 1000 вузлів із $k=1$ медіанний час прибуття до останньої спільноти становив 2785 у \Full{} і 2877 у \TS{}, а $\eW=3.34\%$ (рисунок~\ref{fig:chain_cliques}). Пермутаційний поріг для цієї комірки становив 5.1\% ($p=0.32$), тому за цього обсягу вибірки тест не виявив відмінності між двома термінальними розподілами. Функції розподілу перетинаються, отже різниця не є сталим запізненням. Кліка з 1000 вузлів у середньому насичується за $2H_{999}/(\beta n)\approx5$ одиниць часу, тоді як середній час переходу через ребро дорівнює 333 ($\rho=0.002$). Тому агрегована інтенсивність відрізняється від точної лише протягом незначної частини кожного переходу.

\begin{figure}[tbp]
\centering
\includegraphics[width=\linewidth]{Images/twoscale/corrected_chain_20261008/chain_cliques_n1000_k1_propagation_cdf_random_initial.pdf}
\caption{Ланцюг із десяти клік по 1000 вузлів з одним ребром між послідовними кліками ($k=1$, $\beta=0.003$, 400 запусків для кожної моделі). Ліворуч: поточкова медіана кількості інфікованих вузлів із міжквартильними смугами, які показують розкид між реалізаціями, а не довірчі інтервали. Праворуч: емпірична функція розподілу часу прибуття до десятої спільноти ($T_9$). Чотири значення прибуття \Full{} і два \TS{} лежать поза показаним діапазоном; усі вони спостережені.}
\label{fig:chain_cliques}
\end{figure}

\subsection{Ланцюги випадкових спільнот}
\label{sec:res_chain_random}

\emph{Розріджені спільноти з довгими шляхами дають систематично ранні прибуття, відносна величина яких змінюється значно менше, ніж абсолютна, тоді як абсолютне випередження накопичується.} На рисунку~\ref{fig:chain_heatmap} показано кумулятивні часи прибуття у ланцюгах спільнот із 1000 вузлів. Для цільового ASPL 2--3 термінальне $\eW$ становило 1.3--2.8\%, що близько до його вибіркового порога. За ASPL 6 модель \TS{} досягала останньої спільноти раніше на 13.8\% за медіаною для $k=1$ і на 10.1\% для $k=5$ ($\eW=13.6\%$ і 10.8\%). У рядку ASPL 5 за $k=1$ відносне медіанне випередження становило 9.6\% у першій спільноті та 8.7\% в останній, тоді як його абсолютна величина зросла приблизно зі 107 до 896 одиниць часу. Отже, близькі відносні похибки вздовж ланцюга можуть відповідати дедалі більшому абсолютному випередженню у віддалених спільнотах. У першому стовпці рядка клік зсув медіани додатний ($+12.6\%$ і $+11.5\%$). Окремо, для $k=1$ зсув середнього першого переходу дорівнює $+4.9\%$ із 95\% інтервалом $[-8.0,18.3]\%$, який містить нуль. Цей інтервал стосується середнього, а не медіани. Для одиничних переходів також характерний вищий вибірковий поріг відстаней, ніж для кумулятивних.

\begin{figure}[tbp]
\centering
\includegraphics[width=\linewidth]{Images/twoscale/corrected_chain_20261008/aspl_chain_median_arrival_error_map_labeled.pdf}
\caption{Точність кумулятивних часів прибуття у ланцюгах із десяти спільнот, кожна з яких має $n=1000$ вузлів, залежно від цільового ASPL (рядки) та досягнутої спільноти ($c=1,\dots,9$, стовпці), для $k=1$ (угорі) та $k=5$ (унизу). Ліворуч: відносний зсув медіани $b_c$ (від'ємний: \TS{} прибуває раніше). Праворуч: відносна відстань Вассерштайна. Для кожної комірки використано 400 парних запусків; невизначеність не показано (термінальні прибуття див. на рисунку~\ref{fig:accuracy_map}).}
\label{fig:chain_heatmap}
\end{figure}

\subsection{Розмір спільноти та ширина інтерфейсу}
\label{sec:res_size}

\emph{Та сама довжина внутрішнього шляху є значно менш сприятливою в малих спільнотах.} На рисунку~\ref{fig:accuracy_map} наведено термінальну розбіжність для всіх 48 конфігурацій серії за розміром. 95-й процентиль відстані за внутрішньопарних перестановок був у межах від 1.6\% до 5.9\%. У 16 конфігураціях з ASPL $\le2$ значення $\eW$ становило 1.1--4.5\%, а в 14 із них перестановочний тест не виявив відмінності між розподілами ($p\ge0.05$). За окремими тестами без поправки на множинні порівняння, у 31 із 32 конфігурацій з ASPL $\ge3$ перестановочний тест виявив відмінність між розподілами ($p<0.01$), а зсув медіани був від'ємним із блоковим бутстреп-інтервалом, що не містив нуль. Винятком була конфігурація $n=1000$, ASPL 3, $k=1$, для якої $\eW=1.6\%$. За фіксованого ASPL $\ge3$ розбіжність монотонно зменшувалася зі збільшенням розміру спільноти для обох ширин інтерфейсу. За ASPL 6 значення $\eW$ зменшилося з 40.9\% для $n=50$ до 13.6\% для $n=1000$ за $k=1$, а за $k=5$ --- з 32.6\% до 10.8\%. Ширші інтерфейси зазвичай зменшували розбіжність у розріджених спільнотах, але не усували її. Частка переходів \TS{}, що передували зараженню будь-якого вузла-моста у спільноті-джерелі, зросла з 0.1--2.3\% за ASPL 1 до 31--51\% за ASPL 6.

\begin{figure}[tbp]
\centering
\includegraphics[width=\linewidth]{Images/twoscale/corrected_chain_20261008/fig_accuracy_map.pdf}
\caption{Точність термінального прибуття у ланцюгах із десяти спільнот для всіх 48 конфігурацій серії за розміром. Колір і верхні числа: відносна відстань Вассерштайна $\eW$ для $T_9$ (\%); напівжирні числа позначають комірки з $p<0.05$. Числа курсивом: відносний зсув медіани $b_9$ (\%; від'ємне значення означає, що \TS{} прибуває раніше). Заштриховані комірки: перестановочний тест не виявив відмінності між розподілами ($p\ge0.05$). Значення $p$ наведено без поправки на множинні порівняння. Для кожної комірки використано 400 парних запусків.}
\label{fig:accuracy_map}
\end{figure}

\subsection{Зв’язок похибки зі співвідношенням часових масштабів}
\label{sec:res_descriptor}

\emph{У цій серії показник $\rho=\beta\sigma_T$ має вищу спостережену рангову кореляцію з величиною похибки, ніж ASPL.} Для кожної конфігурації ми обчислювали $\sigma_T$ як стандартне відхилення часів зараження вузлів у кожній спільноті-джерелі запусків \Full{}, усереднене за спільнотами й запусками. На рисунку~\ref{fig:descriptor}(a,b) термінальні $\eW$ і $b_9$ зображено залежно від $\rho$. Рангова кореляція між $\eW$ і $\rho$ становила 0.91 для всіх 48 конфігурацій і 0.98 для 32 конфігурацій з ASPL $\ge3$; відповідні значення для цільового ASPL становили 0.82 і 0.66. У 18 конфігураціях із $\rho<0.25$ виконувалося $\eW\le4.5\%$, а для 15 із них перестановочний тест не виявив відмінності між розподілами. Для всіх 30 конфігурацій із $\rho\ge0.25$ тест виявив відмінності ($p<0.05$). Цей зв’язок є описом отриманих даних, а не універсальним законом похибки. Для $\rho\ge0.3$ апроксимація в логарифмічних координатах дала $\eW\approx15\rho\,\%$ (нахил 1.01; медіанне відносне відхилення від $15\rho$, нормоване на спостережене $\eW$, становило 11\%, максимальне 39\%) і $b_9\approx-\eW$, тобто розбіжність майже повністю була раннім зсувом. Порогові значення та наближення $15\rho$ описують ту саму вибірку, за якою їх отримано. На незалежних мережах їхню здатність прогнозувати похибку не перевірено.

\begin{figure}[tbp]
\centering
\includegraphics[width=\linewidth]{Images/twoscale/corrected_chain_20261008/fig_descriptor.pdf}
\caption{(a) Відносна відстань Вассерштайна та (b) відносний зсув медіани термінального прибуття залежно від $\rho=\beta\sigma_T$ для всіх 48 конфігурацій ланцюга (кольори: розмір спільноти; зафарбовані кола: $k=1$; порожні трикутники: $k=5$; лінії сполучають ASPL 1--6 за фіксованих $n$ і $k$). Сірі позначки в (a): 95-й процентиль внутрішньопарного пермутаційного порога кожної конфігурації. (c) Локальний середній зсув першого переходу порівняно зі зсувом переходів 2--9 для конфігурацій з ASPL $\ge3$, із 95\% блоковими бутстреп-інтервалами зсуву першого переходу.}
\label{fig:descriptor}
\end{figure}

Рисунок~\ref{fig:descriptor}(c) розділяє окремі переходи. Згідно із зауваженням~\ref{rem:placement}, у нашому дизайні перший перехід залежить лише від механізмів на боці джерела та чисельних механізмів, тоді як наступні переходи також залежать від розміщення після першого імпорту й надлишкових імпортів. Для $k=1$ наступні переходи мали сильніше зміщення, ніж перший, у 15 із 16 конфігурацій з ASPL $\ge3$. За окремими блоковими інтервалами без поправки на множинні порівняння різниця була відмінною від нуля у 9 з них, зокрема для всіх чотирьох розмірів за ASPL 6; наприклад, $-27.5\%$ для першого переходу порівняно з $-42.5\%$ для переходів 2--9 за $n=50$. Для $k=5$ різниця була меншою (14 із 16 точкових оцінок від'ємні, 4 статистично відмінні від нуля). Така залежність від $k$ сумісна з поясненням через вичерпання фізичних каналів інтерфейсу: одне фізичне ребро допускає один перехід, тоді як п'ять ребер допускають до п'яти, тому надлишкові агреговані імпорти мають менший вплив. Оскільки між конфігураціями використано спільне випадкове засівання, ми вважаємо цю закономірність описовим, а не остаточним свідченням. Саме порівняння часів переходів не розділяє внески розміщення імпорту, повторних імпортів і чисельної синхронізації. Для цього потрібен окремий аналіз типів подій і контрольні варіанти з точними інтенсивностями або розміщенням імпорту лише на справжніх вузлах-мостах; у наведеному порівнянні їх не використано.
